\SourceLeaf{412}%
\par
Damit gelten die folgenden Beziehungen
\[
\left.\begin{aligned}
tu.\sin p&=2a\\ tv.\sin q&=2x\\ tw.\sin r&=2e\\
vw.\sin(r-q)&=2d\\ uw.\sin(r-p)&=2z\\ uv.\sin(q-p)&=2y
\end{aligned}\right\}\quad
\begin{gathered}\text{daraus folgt mithilfe des Lemmas der}\\
\textit{Th. M. C. C.}\\ ad-xz+ey=0\qquad\text{(I)}\end{gathered}
\]
\begingroup\small
\[
\left.\begin{aligned}b+d+x&=\omega\\a+d+y&=\omega\\a+c+z&=\omega\end{aligned}\right\}\quad
\begin{gathered}\text{setzt man die Werte für }x,y,z\text{ daraus in (I) ein,}\\\text{so erhält man}\end{gathered}
\]
\endgroup
\[ad-(\omega-b-d)(\omega-a-c)+e(\omega-a-d)=0\]
oder ausmultipliziert
\[\omega\omega-(a+b+c+d+e)\omega+(ab+bc+cd+de+ea)=0\]
\begin{flushright}\textsc{Schumacher}.\}\end{flushright}
